\documentclass{article}
\usepackage{amsmath,amssymb}
\usepackage{xurl}
\usepackage[hidelinks]{hyperref}
% Shared commands for notes in docs/paper-gaps/.

\DeclareUnicodeCharacter{2080}{\ifmmode{}_{0}\else\textsubscript{0}\fi}
\DeclareUnicodeCharacter{2081}{\ifmmode{}_{1}\else\textsubscript{1}\fi}
\DeclareUnicodeCharacter{2082}{\ifmmode{}_{2}\else\textsubscript{2}\fi}
\DeclareUnicodeCharacter{2099}{\ifmmode{}_{n}\else\textsubscript{n}\fi}

\newcommand{\Lean}{\textsc{Lean}}
\ExplSyntaxOn
\NewDocumentCommand{\leanid}{m}{
  \ifmmode
    \textsf{\detokenize{#1}}
  \else
    \group_begin:
    \urlstyle{sf}
    \tl_set:Nn \l_tmpa_tl {#1}
    \tl_replace_all:Nnn \l_tmpa_tl {\_} {_}
    \exp_args:NV \path \l_tmpa_tl
    \group_end:
  \fi
}
\ExplSyntaxOff
\providecommand{\tr}{\operatorname{tr}}
\newcommand{\tnleanIssueBase}{https://github.com/LionSR/TNLean/issues/}
\newcommand{\tnleanPrBase}{https://github.com/LionSR/TNLean/pull/}
\newcommand{\tnleanDocBase}{https://sirui-lu.com/TNLean/docs/}
\newcommand{\ghissue}[1]{\href{\tnleanIssueBase#1}{GitHub issue \##1}}
\newcommand{\ghpr}[1]{\href{\tnleanPrBase#1}{PR \##1}}
\newcommand{\doclink}[1]{\href{\tnleanDocBase#1.html}{\path{#1}}}

% Dirac notation and the identity operator, as in blueprint/src/macros/common.tex.
% \providecommand keeps note-local definitions authoritative.
\usepackage{dsfont}
\providecommand{\ket}[1]{|#1\rangle}
\providecommand{\bra}[1]{\langle #1|}
\providecommand{\braket}[2]{\langle #1 | #2 \rangle}
\providecommand{\ketbra}[2]{|#1\rangle\!\langle #2|}
\providecommand{\Id}{\mathds{1}}
\providecommand{\one}{\mathds{1}}
\providecommand{\C}{\mathbb{C}}
\providecommand{\MN}[1]{M_{#1}(\C)}
\providecommand{\MD}{M_D(\C)}

% External referencing for paper sources.  The build helper writes sanitized
% auxiliary files under build/paper-aux/ and excludes duplicated source labels.
\usepackage{xr}

% Import a generated paper auxiliary file with a caller-supplied label prefix.
% The candidate paths support builds from docs/paper-gaps/, the repository root,
% and build/paper-aux/ itself.
\newcommand{\importpaperaux}[2]{%
  \IfFileExists{../../build/paper-aux/#2.aux}{%
    \externaldocument[#1]{../../build/paper-aux/#2}%
  }{%
    \IfFileExists{build/paper-aux/#2.aux}{%
      \externaldocument[#1]{build/paper-aux/#2}%
    }{%
      \IfFileExists{#2.aux}{%
        \externaldocument[#1]{#2}%
      }{}%
    }%
  }%
}

% General external reference macros
\newcommand{\paperref}[2]{\ref{#1-#2}}
\newcommand{\papereqref}[2]{(\ref{#1-#2})}

% CPSV16 source (1606.00608 / MPDO-22-12-17-2.tex) external document and shortcuts
\importpaperaux{cpsv16-}{1606.00608/MPDO-22-12-17-2-tnlean}
\newcommand{\cpsvref}[1]{\paperref{cpsv16}{#1}}
\newcommand{\cpsveqref}[1]{\papereqref{cpsv16}{#1}}

% Verdict marker: \gapnote{<kind>}{<status>}. Kind is the policy
% classification (clarification, local-correction, scope-restriction,
% unfaithful, false-source, open-gap); status is open, wip (elimination
% underway), resolved, or historical. Machine-read by the site index;
% prints a verdict line.
\newcommand{\gapnote}[2]{\par\noindent\textbf{Verdict:} #1 (#2).\par}

\title{Sequential Factorization of the Blocked Isometry: Positive Block Length}
\author{}
\date{2026-09-26}
\begin{document}
\maketitle
\gapnote{scope-restriction}{open}

\section{What the source states}
Let $A$ be a tensor of physical dimension $d$ and bond dimension $D$, and let
$\mathcal B_q\colon\C^D\otimes\C^D\to(\C^d)^{\otimes q}$ be its map blocked
over $q$ sites, with polar decomposition $\mathcal B_q=VP$. Equations~(13)--(15)
of~\cite{Malz2024Preparation} write the isometry as an ordered product of $q$
isometries acting on one site and a bond each,
\begin{align}
    V
     & =V_q\cdots V_1,
    \label{eq:mswc24_sequential_product}
\end{align}
where $V_k\colon\C^{D'_k}\to\C^{D'_{k+1}}\otimes\C^d$, $D'_1=D^2$,
$D'_{q+1}=1$ and $D'_k\leq D^2$. The source states no lower bound on the block
length $q$; the construction factors the blocked map site by site and blocks
at least one site.

\section{The empty block}
For $q=0$ the blocked map has a single matrix, the empty product
$\Id_D$, and the output space $(\C^d)^{\otimes0}$ is $\C$. The one matrix
$\Id_D$ spans the full matrix algebra only when $D\leq1$. When $D=1$,
$V=1$ and (\ref{eq:mswc24_sequential_product}) is the empty product with
$D'_1=D'_{q+1}=1$, so the statement holds trivially. When $D=0$ the space
$\C^D\otimes\C^D$ is zero, and the requirement $D'_1=D^2=0$ contradicts
$D'_{q+1}=1$ at $q=0$, where these are the same bond. The empty block
carries no content of the preparation scheme, which blocks at least one site.

\section{The restriction adopted}
The sequential factorization is stated for block lengths $q\geq1$ and injective
$\mathcal B_q$.\footnote{Formal declaration:
    \leanid{MPSPreparation.exists_isometric_chain_polarIsoMatrix} in
    \doclink{TNLean/MPS/Preparation/SequentialFactorization}.}
For every $q\geq1$ the conclusion is exactly
(\ref{eq:mswc24_sequential_product}) in matrix product form, with no further
hypothesis. Elimination: replace $q\geq1$ by a positive bond dimension $D\geq1$,
which the source takes for granted, and treat the case $q=0$, $D=1$ directly
as the empty product.

\bibliographystyle{alpha}
\bibliography{references}
\end{document}
